\chapter{Regresi dan Statistika Terapan}

Sebuah kafe kecil mencatat sembilan hari dagang. Dua angka ditulis setiap malam:
banyak pesanan daring dalam puluhan, dan jam lembur dapur. Pemilik kafe membutuhkan
satu angka untuk satu pertanyaan: bila pesanan bertambah sepuluh, berapa tambahan jam
lembur yang harus ia siapkan besok?

\begin{center}
\small
\begin{tabular}{l*{9}{r}}
\hline
hari               & 1 & 2 & 3 & 4 & 5 & 6 & 7 & 8 & 9 \\
pesanan (puluhan)  & 1 & 2 & 3 & 4 & 5 & 6 & 7 & 8 & 9 \\
lembur (jam)       & 10 & 12 & 13 & 14 & 16 & 18 & 19 & 20 & 22 \\
\hline
\end{tabular}
\end{center}

Jawaban cepat tersedia tanpa alat apa pun: dari hari pertama ke hari terakhir pesanan
naik delapan puluhan dan lembur naik dua belas jam, jadi kira-kira $1{,}5$ jam per
sepuluh pesanan. Kuadrat terkecil, yang akan kita susun di bab ini, memberi
$1{,}466667$. Kedua angka itu dekat, dan justru karena dekat kita tergoda menganggap
persoalannya selesai. Tiga pertanyaan penting tidak terjawab oleh keduanya. Seberapa
jauh angka itu bergeser bila sembilan hari itu diganti sembilan hari lain? Bila besok
pesanannya delapan puluh, berapa jam yang harus disiapkan? Dan bila satu hari ternyata
salah catat, berapa besar kerusakannya?

Bab 10 sudah memberi gambarannya: kuadrat terkecil adalah proyeksi ortogonal vektor
pengamatan ke ruang kolom matriks rancangan, dan persamaan normal menyatakan bahwa
sisaannya tegak lurus ruang itu. Yang belum ada di sana adalah ketidakpastian. Bab 20
sudah memberi selang kepercayaan dan uji hipotesis, tetapi untuk satu parameter dari
pengamatan yang saling bebas dan berbagi satu sebaran. Di sini parameter yang dicari
adalah kemiringan, nisbah dua jumlah yang keduanya acak, dan pengamatan tidak lagi
sederajat: titik di ujung rentang lebih berpengaruh daripada titik di tengah. Bab ini
menyambung keduanya dan berhenti di tempat yang jujur: seberapa jauh kesimpulan boleh
dibawa.

\begin{untukanak}
Seorang nenek menjual es kelapa di depan rumah. Ia tidak pernah mencatat apa pun
dengan rapi, tetapi ingatannya cukup: hari yang panas, dagangannya habis sebelum
sore. Hari yang mendung, separuh es mencair di termos. Suatu hari cucunya datang
dengan buku tulis dan mencatat Sembilan hari berturut-turut: berapa kira-kira panasnya,
dan berapa gelas yang terjual.

Dari sembilan catatan itu cucunya menarik satu garis lurus: makin panas, makin banyak
yang terjual. Neneknya mengangguk, tetapi lalu bertanya, ``kalau besok lebih panas
sedikit, berapa gelas yang harus saya siapkan?'' Cucunya terdiam. Garis itu mengatakan
arahnya, bukan angkanya. Lagi pula, satu hari pernah hujan deras padahal panas, dan
penjualannya jelek sekali; titik itu jauh dari garis.

Di bab ini kita akan tahu mengapa garis yang paling pas tetap bukan ramalan yang pasti,
mengapa satu hari yang aneh dapat menggeser garis lebih jauh daripada lima hari biasa,
dan mengapa orang yang menjual angka harus menyebutkan seberapa lebar ketidakpastiannya.
\end{untukanak}

\section{Sembilan hari dan satu garis}
\tingkat{dasar}

Sembilan pasang angka, dua bilangan yang dicari, sembilan persamaan. Sistem seperti ini
sudah kita temui di Bab 8: lebih banyak persamaan daripada bilangan yang tidak diketahui,
sehingga solusi tepat hanya ada bila kesembilan titik itu kebetulan segaris. Tidak ada
garis yang lewat persis di semua titik di atas, dan itu bukan kegagalan pengukuran,
melainkan keadaan biasa. Yang bisa dicari hanyalah garis yang paling dekat, dan
``paling dekat'' harus didefinisikan lebih dahulu.

Kalimat sehari-harinya: kita mencari satu garis yang jumlah kesalahannya sekecil mungkin,
dengan kesalahan dihitung sebagai jarak tegak dari titik ke garis.

\kamus{regresi}{regression}{menarik satu garis yang paling dekat ke semua titik, bukan
garis yang lewat di semua titik}

\begin{definisi}[Model regresi linear sederhana]
\index{regresi}
Untuk pengamatan $i=1,\dots,n$,
\[
y_i = a + b\,x_i + \varepsilon_i ,
\]
dengan $x_i$ peubah penjelas yang teramati, $y_i$ peubah yang dijelaskan, $a$ titik
potong, $b$ kemiringan, dan $\varepsilon_i$ galat yang tidak teramati. Model disebut
\emph{linear} karena linear dalam parameter $a$ dan $b$, bukan karena garisnya lurus
dalam $x$.
\end{definisi}

Satu perbedaan kecil yang menentukan seluruh bab: $\varepsilon_i$ dan sisaan bukan
benda yang sama. Galat $\varepsilon_i$ adalah selisih antara nilai sebenarnya dan garis
populasi yang tidak pernah kita lihat; sisaan $e_i=y_i-\hat y_i$ adalah selisih antara
nilai teramati dan garis yang kita pasang dari data itu sendiri. Sisaan dapat dihitung
dan digambar, galat hanya dapat diasumsikan. Setiap kali sebuah buku menyebut asumsi
galat lalu memeriksanya lewat gambar sisaan, ada satu lompatan yang tidak boleh
dilupakan: yang diperiksa adalah pengganti, bukan barang aslinya.

\section{Kuadrat terkecil sebagai proyeksi}
\tingkat{menengah}

Susun $y\in\mathbb{R}^n$ sebagai vektor pengamatan dan $X$ matriks $n\times2$ yang
kolomnya satu dan $x$. Mencari garis terbaik berarti memilih $\hat y$ di ruang kolom $X$
yang paling dekat ke $y$: persoalan Bab 10, dengan tambahan bahwa kita butuh bentuk
eksplisit koefisiennya untuk menghitung ketidakpastiannya.

\begin{definisi}[Penaksir kuadrat terkecil]
\index{kuadrat terkecil}
$\hat\beta=(\hat a,\hat b)$ adalah penaksir kuadrat terkecil bila ia meminimumkan
\[
\operatorname{RSS}(\beta)=\|y-X\beta\|^2=\sum_{i=1}^{n}\bigl(y_i-a-bx_i\bigr)^2 .
\]
\end{definisi}

\kamus{kuadrat terkecil}{least squares}{menjumlahkan semua kesalahan setelah
dikuadratkan, lalu memilih garis yang jumlahnya paling kecil}

\begin{teorema}[Persamaan normal]
\index{persamaan normal}
$\hat\beta$ meminimumkan RSS jika dan hanya jika
\[
X^{\top}(y-X\hat\beta)=0 ,
\qquad\text{yaitu}\qquad X^{\top}X\hat\beta=X^{\top}y .
\]
Bila kolom $X$ bebas linear, $\hat\beta=(X^{\top}X)^{-1}X^{\top}y$ tunggal.
\end{teorema}

\begin{proof}[Bukti]
Tulis $f(\beta)=\|y-X\beta\|^2=y^{\top}y-2\beta^{\top}X^{\top}y+\beta^{\top}X^{\top}X\beta$.
Gradiennya $-2X^{\top}y+2X^{\top}X\beta$, jadi syarat stasioner tepat persamaan normal.
Karena $X^{\top}X$ definit positif bila kolomnya bebas linear, $f$ konveks ketat dan
titik stasioner itu minimum global.

Arah sebaliknya lebih informatif. Andaikan $X^{\top}(y-X\hat\beta)=0$. Untuk sebarang
$\beta$,
\[
\|y-X\beta\|^2=\|y-X\hat\beta\|^2+\|X(\hat\beta-\beta)\|^2
+2(\hat\beta-\beta)^{\top}X^{\top}(y-X\hat\beta),
\]
dan suku silang lenyap. Jadi $\operatorname{RSS}(\beta)\ge\operatorname{RSS}(\hat\beta)$
untuk semua $\beta$, dengan kesamaan hanya bila $X\beta=X\hat\beta$.
\end{proof}

Pembuktian itu memuat seluruh isi geometrisnya: sisaan tegak lurus setiap kolom $X$.
Karena kolom pertama vektor satu dan kolom kedua $x$, berlaku $\sum_i e_i=0$ dan
$\sum_i x_ie_i=0$, dan kedua kesamaan itulah yang menopang dekomposisi jumlah kuadrat
di bagian berikutnya.

\begin{proposisi}[Matriks topi]
\index{matriks topi}
Dengan $H=X(X^{\top}X)^{-1}X^{\top}$, berlaku $\hat y=Hy$, $H^{\top}=H$, $H^2=H$,
$\operatorname{tr}H=2$, dan $0\le h_{ii}\le1$ dengan $\sum_i h_{ii}=2$. Nilai $h_{ii}$
disebut \emph{leverage} titik $i$.
\end{proposisi}

\begin{proof}[Bukti]
Simetri dan idempotensi diperoleh langsung dari definisi:
\[
H^2=X(X^{\top}X)^{-1}X^{\top}X(X^{\top}X)^{-1}X^{\top}=H .
\]
Karena itu $H$ adalah matriks proyeksi ortogonal,
sehingga $h_{ii}=e_i^{\top}He_i=\|He_i\|^2\ge0$ dan $h_{ii}\le\|e_i\|^2=1$. Terakhir
$\operatorname{tr}H=\operatorname{tr}\bigl((X^{\top}X)^{-1}X^{\top}X\bigr)
=\operatorname{tr}I_2=2$.
\end{proof}

Sekarang angkanya, karena seluruh bab bergantung pada satu perhitungan yang dapat
diperiksa ulang. Untuk data kafe,
\[
X^{\top}X=\begin{pmatrix}9&45\\45&285\end{pmatrix},
\qquad X^{\top}y=\begin{pmatrix}144\\808\end{pmatrix},
\qquad \det(X^{\top}X)=9\cdot285-45^2=540 .
\]
Maka
\[
\hat\beta=\frac1{540}\begin{pmatrix}285&-45\\-45&9\end{pmatrix}
\begin{pmatrix}144\\808\end{pmatrix}
=\frac1{540}\begin{pmatrix}4680\\792\end{pmatrix}
=\begin{pmatrix}8{,}666667\\1{,}466667\end{pmatrix}.
\]
Angka $808$ layak diperiksa dengan tangan: $\sum x_iy_i=10+24+39+56+80+108+133+160+198$.
Angka $4680$ juga: $285\cdot144-45\cdot808=41040-36360$.

Satu teorema lagi diperlukan sebelum kita boleh bicara tentang mutu penaksir. Teorema
itu menjelaskan mengapa kuadrat terkecil dipakai di hampir semua bidang, dan sekaligus
mengapa ia bukan pilihan suci.

\begin{teorema}[Gauss--Markov]
\index{gauss-markov}
Andaikan (A1) $y=X\beta+\varepsilon$; (A2) $\mathbb{E}[\varepsilon\mid X]=0$;
(A3) $\operatorname{Var}(\varepsilon\mid X)=\sigma^2I$; (A4) kolom $X$ bebas linear.
Maka $\hat\beta$ tak bias dengan $\operatorname{Var}(\hat\beta)=\sigma^2(X^{\top}X)^{-1}$,
dan di antara semua penaksir linear tak bias, $\hat\beta$ yang bervariansi terkecil.
Penaksir seperti itu disebut \emph{BLUE} (\emph{best linear unbiased estimator}).
\end{teorema}

\begin{proof}[Bukti]
Ketakbiasan: $\mathbb{E}[\hat\beta\mid X]=(X^{\top}X)^{-1}X^{\top}X\beta=\beta$ memakai
(A2). Variansi: $\hat\beta-\beta=(X^{\top}X)^{-1}X^{\top}\varepsilon$, jadi
$\operatorname{Var}(\hat\beta)=(X^{\top}X)^{-1}X^{\top}(\sigma^2I)X(X^{\top}X)^{-1}
=\sigma^2(X^{\top}X)^{-1}$.

Bagian ketiga tidak dibuktikan penuh di sini. Sketsanya: setiap penaksir linear tak
bias berbentuk $\tilde\beta=Cy$ dengan $CX=I$, sehingga $\tilde\beta-\beta=C\varepsilon$
dan $\operatorname{Var}(\tilde\beta)=\sigma^2CC^{\top}$. Menulis $C=(X^{\top}X)^{-1}X^{\top}+D$
memberi $DX=0$ dan
\[
CC^{\top}-(X^{\top}X)^{-1}=DD^{\top},
\]
yang matriks semi-definit positif, sehingga setiap komponen $\tilde\beta$ bervariansi
tak lebih kecil daripada komponen yang sepadan dari $\hat\beta$. Bukti lengkap ada di
Seber dan Lee, \emph{Linear Regression Analysis}, Bab 3, dan di Greene,
\emph{Econometric Analysis}, Bab 3.
\end{proof}

\section{Koefisien dan satuannya}
\tingkat{dasar}

Kemiringan $1{,}466667$ tidak berarti apa pun sebelum satuannya disebut. Satuannya jam
per puluhan pesanan: setiap tambahan sepuluh pesanan menaikkan ramalan lembur sebesar
$1{,}466667$ jam, yaitu $88$ menit, atau $8{,}8$ menit per pesanan. Kalimat itu dapat
dipakai pemilik kafe; kalimat ``koefisiennya $1{,}47$'' tidak.

Titik potong $8{,}666667$ jam adalah ramalan untuk $x=0$, yaitu hari tanpa satu pun
pesanan daring. Tidak ada hari seperti itu dalam data, dan di luar rentang $1$ sampai
$9$ garisnya tidak dijamin berlaku. Menyebut angka itu ``biaya tetap lembur'' adalah
\index{ekstrapolasi}ekstrapolasi yang dibungkus nama: kita memakai garis di luar
tempat ia dipasang. Yang aman: titik potong adalah bilangan yang membuat garis lewat
$(\bar x,\bar y)$, sebab $\hat a=\bar y-\hat b\bar x$; ia berguna untuk perhitungan,
bukan untuk tafsiran.

Perubahan satuan tidak mengubah kesimpulan. Bila $x$ ditulis dalam pesanan, bukan
puluhan, kemiringannya $0{,}146667$ jam per pesanan dan galat bakunya $0{,}004714$,
keduanya terbagi sepuluh; $t$ dan $R^2$ tidak bergerak. Mengalikan satu peubah dengan
konstanta $c$ mengalikan koefisien dan galat bakunya dengan $c$ atau $1/c$, sedangkan
nisbah keduanya tetap, sehingga laporan yang menyebut $t$ tanpa satuan tidak terbaca.

\begin{tangan}
Siapkan kertas berpetak, sembilan koin, dan seutas benang sepanjang sisi kertas. Letakkan
koin pada petak sesuai sembilan pasang angka di awal bab: mendatar pesanan, tegak lembur,
satu petak satu satuan. Ikat ujung benang pada dua ujung kertas sehingga benang dapat
digeser sebagai garis lurus. Geser dan putar benang sampai jumlah jarak tegak semua koin
ke benang terasa paling kecil; jangan mengukur, cukup bandingkan dengan mata, lalu
kunci dengan selotip.

Yang harus terlihat: dua orang akan menaruh benang di tempat yang sedikit berbeda, dan
selisihnya paling besar ketika engkau menggeser ujung benang. Sekarang angkat koin
paling kiri satu petak ke atas, dan perhatikan betapa besar benang harus berputar.
Kembalikan, lalu angkat koin tengah satu petak: hampir tidak ada yang berubah. Itulah
leverage, dan engkau baru saja melihatnya dengan tangan sebelum bertemu rumusnya.
\end{tangan}

\section{R kuadrat dan korelasinya}
\tingkat{menengah}

Sesudah garisnya ada, orang bertanya ``seberapa bagus''. Jawaban yang paling sering
dipakai adalah satu angka yang membandingkan sisaan dengan sebaran asli $y$. Kalimat
sehari-harinya: berapa bagian dari goyangan $y$ yang berhasil dijelaskan garis, dan
berapa bagian yang tertinggal.

\begin{definisi}[$R^2$]
\index{koefisien determinasi}
Dengan $\mathrm{SST}=\sum_i(y_i-\bar y)^2$, $\mathrm{SSE}=\sum_i e_i^2$, dan
$\mathrm{SSR}=\sum_i(\hat y_i-\bar y)^2$, koefisien determinasi adalah
\[
R^2=1-\frac{\mathrm{SSE}}{\mathrm{SST}}=\frac{\mathrm{SSR}}{\mathrm{SST}} .
\]
\end{definisi}

\kamus{koefisien determinasi}{coefficient of determination}{bagian dari goyangan yang
berhasil dijelaskan garis, dinyatakan sebagai pecahan dari keseluruhan}

\begin{teorema}[Dekomposisi jumlah kuadrat]
Bila $X$ memuat kolom satu, maka $\mathrm{SST}=\mathrm{SSR}+\mathrm{SSE}$ dan
$0\le R^2\le1$.
\end{teorema}

\begin{proof}[Bukti]
Tulis $y-\bar y\mathbf{1}=(\hat y-\bar y\mathbf{1})+e$. Karena $\sum e_i=0$, vektor $e$
tegak lurus $\bar y\mathbf{1}$, dan menurut persamaan normal $e$ juga tegak lurus
$\hat y$. Jadi kedua suku di ruas kanan saling tegak lurus, dan teorema Pythagoras
memberi $\|y-\bar y\mathbf{1}\|^2=\|\hat y-\bar y\mathbf{1}\|^2+\|e\|^2$.
\end{proof}

Untuk data kafe, $\mathrm{SST}=130$, $\mathrm{SSE}=0{,}933333$, sehingga
$R^2=1-0{,}933333/130=0{,}992821$. Pada regresi sederhana, $R^2$ sama dengan kuadrat
korelasi: $r=88/\sqrt{60\cdot130}=0{,}996404$ dan $r^2=0{,}992821$. Kesamaan itu
khusus untuk satu peubah penjelas; begitu penjelasnya dua, $R^2$ bukan lagi kuadrat
korelasi apa pun.

Angka $0{,}992821$ mudah membuat orang berhenti berpikir, jadi mari kita jinakkan. Yang
dijelaskan adalah $99{,}28$ persen \emph{goyangan} lembur, dan goyangan itu diukur
terhadap rata-ratanya sendiri. Sisaannya masih bervariansi $s^2=0{,}933333/7=0{,}133333$,
yaitu simpangan baku $s=0{,}365148$ jam atau $21{,}9$ menit. Pemilik kafe yang menyiapkan
lembur atas dasar garis ini masih bisa meleset lebih dari dua puluh menit. $R^2$ bersifat
nisbi; $s$ bersifat mutlak, dan yang dirasakan orang adalah $s$.

\begin{untukanak}
Bayangkan dua siswa. Yang pertama ikut lomba lari seratus meter dan finis $0{,}2$
detik lebih lambat daripada juara. Yang kedua ikut lomba maraton dan finis dua puluh
menit lebih lambat daripada juara. Keduanya kalah, tetapi tidak dalam arti yang sama.
Dua persepuluh detik itu sangat dekat; dua puluh menit itu sangat jauh. Kalau gurunya
hanya menuliskan ``peringkat dua'' untuk keduanya, kita kehilangan seluruh ceritanya.

Angka yang mirip terjadi pada bab ini. Ada angka yang mengukur \emph{bagian} yang
berhasil ditebak, dan ada angka yang mengukur \emph{besar} kesalahan yang masih
tersisa. Yang pertama bisa tampak hampir sempurna sementara yang kedua tetap besar.
Orang yang hanya membaca angka pertama akan menyiapkan terlalu sedikit, lalu bingung
mengapa ramalannya meleset.

Di bab ini kita akan tahu mengapa dua angka harus disebut berpasangan, mengapa angka
bagus tidak sama dengan ramalan yang sempit, dan kapan garis yang rapi sebaiknya
ditinggalkan.
\end{untukanak}

\begin{proposisi}[Menambah peubah tidak pernah menurunkan $R^2$]
Bila kolom $X_2$ ditambahkan pada $X_1$ sehingga ruang kolomnya memuat ruang kolom
$X_1$, maka $\mathrm{SSE}$ tidak bertambah dan $R^2$ tidak berkurang.
\end{proposisi}

\begin{proof}[Bukti]
Himpunan vektor $X_1\beta_1$ termuat dalam himpunan $X_2\beta_2$, jadi minimum jarak
ke $y$ dihimpun atas himpunan yang lebih besar tidak dapat lebih tinggi. Karena itu
$R^2$ tidak dapat turun, meskipun peubah yang ditambahkan tidak berhubungan dengan $y$.
\end{proof}

Akibat praktisnya penting: membandingkan model dengan $R^2$ mentah selalu memihak model
yang lebih gemuk. Koreksi yang lazim adalah
$R^2_{\text{adj}}=1-(1-R^2)(n-1)/(n-p)$ dengan $p$ banyak parameter. Untuk model
berganda di bagian berikutnya, $R^2=0{,}997951$ turun menjadi $0{,}997496$ ketika $p=3$
dan $n=12$; pada $n$ kecil dengan banyak peubah koreksi itu dapat membalik urutan.

\section{Regresi berganda dan kolinearitas}
\tingkat{lanjut}

Sampai sini satu peubah penjelas. Kenyataannya hampir setiap persoalan menuntut lebih
dari satu, dan di situlah muncul penyakit khasnya: dua penjelas yang membawa informasi
hampir sama.

Dua belas cabang sebuah kedai dicatat: penjualan harian $y$ dalam juta rupiah, banyak
karyawan $x_1$, dan luas ruang duduk $x_2$ dalam meter persegi.

\begin{center}
\small
\begin{tabular}{l*{12}{r}}
\hline
cabang   & 1 & 2 & 3 & 4 & 5 & 6 & 7 & 8 & 9 & 10 & 11 & 12 \\
$x_1$    & 6 & 7 & 7 & 8 & 8 & 9 & 9 & 10 & 10 & 11 & 12 & 12 \\
$x_2$    & 18 & 22 & 20 & 26 & 24 & 28 & 26 & 32 & 30 & 34 & 36 & 38 \\
$y$      & 11,2 & 12,4 & 12,0 & 14,1 & 13,6 & 15,2 & 14,8 & 17,0 & 16,4 & 18,1 & 19,4 & 19,9 \\
\hline
\end{tabular}
\end{center}

Kuadrat terkecil memberi
\[
\hat y=2{,}442991+0{,}686916\,x_1+0{,}239252\,x_2,
\]
dengan galat baku $0{,}207583$, $0{,}135366$, dan $0{,}042101$; nilai $t$ berturut-turut
$11{,}7687$, $5{,}0745$, dan $5{,}6828$ pada $9$ derajat kebebasan; $s=0{,}143607$ dan
$R^2=0{,}997951$. Semua tampak sehat. Lalu kita pasang satu penjelas saja: $y$ atas $x_1$
memberi $1{,}446019$, $y$ atas $x_2$ memberi $0{,}450075$, keduanya hampir dua kali
koefisien bergandanya.

Bukan karena salah hitung. Karyawan dan luas ruang hampir saling menentukan: korelasi
keduanya $r_{12}=0{,}986801$. Setelah $x_2$ masuk, $x_1$ hampir tidak membawa
keterangan baru, jadi koefisiennya menyusut menjadi efek tambahannya saja, yaitu
$0{,}687$ juta per karyawan pada luas yang tetap. Koefisien tunggal $1{,}446$ mencampur
efek karyawan dengan efek luas yang biasanya menyertainya, dan jurang antara keduanya
adalah ukuran seberapa jauh kedua peubah itu bertumpang.

\begin{definisi}[Faktor inflasi varians]
\index{varians inflasi}
Untuk peubah penjelas ke-$j$, misalkan $R_j^2$ adalah koefisien determinasi regresi
$x_j$ atas seluruh penjelas lain. Faktor inflasi variansnya
\[
\mathrm{VIF}_j=\frac{1}{1-R_j^2}.
\]
\end{definisi}

\kamus{varians inflasi}{variance inflation factor}{seberapa banyak galat baku sebuah
koefisien digelembungkan oleh penjelas lain yang membawa informasi serupa}

\begin{proposisi}[VIF dua penjelas]
\index{kolinearitas}
Bila hanya ada dua penjelas, $\operatorname{Var}(\hat b_1)=\sigma^2/\bigl(S_{11}
(1-r_{12}^2)\bigr)$, dengan $S_{11}=\sum_i(x_{i1}-\bar x_1)^2$. Jadi galat baku
$\hat b_1$ tergelembung tepat sebesar $\sqrt{\mathrm{VIF}_1}$ dibanding keadaan penjelas
yang saling ortogonal.
\end{proposisi}

\begin{proof}[Bukti]
Untuk dua penjelas, invers matriks $2\times2$ dari bagian terpusat $X^{\top}X$ memberi
unsur pertama $S_{22}/(S_{11}S_{22}-S_{12}^2)$, dan $S_{12}=r_{12}\sqrt{S_{11}S_{22}}$.
Membagi pembilang dan penyebut dengan $S_{22}$ menghasilkan
$\sigma^2/\bigl(S_{11}(1-r_{12}^2)\bigr)$.
\end{proof}

Angkanya dapat diperiksa: $S_{11}=42{,}9167$, $r_{12}^2=0{,}973775$, sehingga
$\mathrm{VIF}=38{,}132232$ dan $\sqrt{\mathrm{VIF}}=6{,}1751$. Galat baku $\hat b_1$
memang $0{,}135366$, sedangkan $\sigma/\sqrt{S_{11}}$ hanya $0{,}021921$; nisbahnya
$6{,}1751$, cocok dengan akar VIF sampai empat angka.

Kolinearitas tidak merusak ramalan; ia merusak tafsiran. Nilai $R^2=0{,}997951$
menunjukkan garisnya masih dekat ke titik, dan ramalan pada $x_1=9$, $x_2=28$ masih
berguna. Yang tidak dapat dipercaya adalah kalimat ``satu karyawan tambahan menaikkan
penjualan $0{,}687$ juta'': selang $95$ persennya $(0{,}380697;\ 0{,}993135)$ lebar,
dan bila data bergeser sedikit koefisien itu dapat berubah banyak tanpa garisnya
berubah banyak.

\section{Selang, uji $t$, dan selang prediksi}
\tingkat{lanjut}

Pertanyaan berikutnya sama dengan Bab 20, hanya bendanya berganti: bukan rata-rata,
melainkan kemiringan. Kita memerlukan satu besaran yang sebarannya diketahui bila
modelnya benar, lalu membaliknya menjadi selang.

Kalimat sehari-harinya: kita mengukur seberapa jauh kemiringan yang teramati dari nol,
dalam satuan goyangan khas kemiringan itu sendiri.

\kamus{selang prediksi}{prediction interval}{selang untuk satu hari yang belum terjadi,
selalu lebih lebar daripada selang untuk rata-rata karena ia harus menanggung kejutan
hari itu juga}

\begin{teorema}[Uji $t$ dan selang untuk satu koefisien]
Andaikan (A1) sampai (A4) dan galatnya normal. Dengan $p$ banyak parameter,
$s^2=\mathrm{SSE}/(n-p)$, dan $\operatorname{se}(\hat b_j)=s\sqrt{[(X^{\top}X)^{-1}]_{jj}}$,
maka
\[
\frac{\hat b_j-\beta_j}{\operatorname{se}(\hat b_j)}\sim t_{n-p},
\]
sehingga selang $1-\alpha$ untuk $\beta_j$ adalah
$\hat b_j\pm t_{n-p,1-\alpha/2}\operatorname{se}(\hat b_j)$, dan uji $H_0\colon\beta_j=0$
menolak pada taraf $\alpha$ bila $|t|>t_{n-p,1-\alpha/2}$.
\end{teorema}

\begin{proof}[Bukti]
$\hat b_j$ adalah kombinasi linear $y$ dengan koefisien tetap, jadi ia normal dengan
harapan $\beta_j$ dan variansi $\sigma^2[(X^{\top}X)^{-1}]_{jj}$ menurut Gauss--Markov.
Selanjutnya $\mathrm{SSE}/\sigma^2\sim\chi^2_{n-p}$ dan bebas dari $\hat\beta$, sebab
sisaan tegak lurus ruang kolom sedangkan $\hat\beta$ seluruhnya di dalam ruang itu.
Nisbah peubah normal baku dengan akar peubah khi-kuadrat dibagi derajat kebebasannya
mengikuti sebaran $t$. Bukti lengkapnya ada di Wasserman, \emph{All of Statistics},
Bab 13.
\end{proof}

Untuk data kafe, $\hat b=1{,}466667$, $s=0{,}365148$, $S_{xx}=60$, jadi
$\operatorname{se}(\hat b)=0{,}047140$. Dengan $t_{7,0{,}975}=2{,}364624$,
\[
\hat b\pm t\,\operatorname{se}(\hat b)=1{,}466667\pm0{,}111470
=(1{,}355197;\ 1{,}578136),
\]
dan $t=31{,}1127$ dengan nilai-p di bawah $10^{-6}$. Menolak $\beta=0$ di sini bukan
prestasi: dengan sisaan sesempit $22$ menit dan rentang $x$ selebar ini, hampir semua
data akan menolaknya. Yang layak dilaporkan justru selangnya, karena ia mengatakan
bahwa tambahan lembur per sepuluh pesanan ada di antara $1{,}36$ dan $1{,}58$ jam.

Sekarang bedakan dua pertanyaan yang sering tertukar. Berapa rata-rata lembur pada hari
dengan delapan puluh pesanan? Dan berapa lembur pada satu hari tertentu dengan delapan
puluh pesanan? Jawaban pertama menyempit sebanding $1/\sqrt n$; jawaban kedua tidak
pernah menyempit di bawah simpangan baku sisaan, sebab satu hari membawa kejutannya
sendiri.

\begin{proposisi}[Dua ragam ramalan]
Untuk $x_0$ tetap dan pengamatan baru $y_0=x_0^{\top}\beta+\varepsilon_0$ yang bebas
dari data lama,
\[
\operatorname{Var}(\hat y_0)=\sigma^2x_0^{\top}(X^{\top}X)^{-1}x_0,
\qquad
\operatorname{Var}(y_0-\hat y_0)=\sigma^2\bigl(1+x_0^{\top}(X^{\top}X)^{-1}x_0\bigr).
\]
\end{proposisi}

\begin{proof}[Bukti]
Yang pertama sudah muncul di Gauss--Markov. Untuk yang kedua, tulis
$y_0-\hat y_0=\varepsilon_0-x_0^{\top}(\hat\beta-\beta)$; kedua suku bebas karena
$\varepsilon_0$ bebas dari data lama, sehingga variansinya jumlah kedua variansi.
\end{proof}

Untuk $x_0=8$ pada data kafe, $x_0^{\top}(X^{\top}X)^{-1}x_0=1/9+9/60=0{,}261111$.
Setengah lebar selang rata-rata $2{,}364624\sqrt{0{,}133333\cdot0{,}261111}=0{,}441209$,
jadi $(19{,}9588;\ 20{,}8412)$ selebar $0{,}882$. Setengah lebar selang prediksi
$2{,}364624\sqrt{0{,}133333\cdot1{,}261111}=0{,}969635$, jadi
$(19{,}4304;\ 21{,}3696)$ selebar $1{,}939$. Nisbah lebarnya
$\sqrt{1{,}261111/0{,}261111}=2{,}198$: selang prediksi lebih dari dua kali lebih lebar,
dan tidak ada jumlah data yang dapat menghapus kelebihan itu. Gambar~\ref{fig:dua-selang}
menunjukkan kedua pita untuk seluruh rentang $x$.

\begin{figure}[htbp]
\centering
\begin{tikzpicture}
\begin{axis}[
  width=11.4cm, height=6.0cm,
  xmin=0.6, xmax=9.4, ymin=8.2, ymax=25.4,
  axis lines=left,
  xtick={1,2,3,4,5,6,7,8,9},
  tick label style={font=\scriptsize},
  xlabel={pesanan daring (puluhan)}, ylabel={jam lembur},
  xlabel style={font=\small}, ylabel style={font=\small},
  legend style={font=\scriptsize,draw=layergrey,at={(0.02,0.98)},
    anchor=north west,row sep=0.5pt},
]
\addplot[accent,very thick,domain=0.7:9.3,samples=60]
  {8.666667+1.466667*x+0.863439*sqrt(1.111111+(x-5)^2/60)};
\addplot[accent,very thick,domain=0.7:9.3,samples=60]
  {8.666667+1.466667*x-0.863439*sqrt(1.111111+(x-5)^2/60)};
\addplot[accent,very thick,domain=0.7:9.3,samples=60,forget plot]
  {8.666667+1.466667*x};
\addplot[layerindigo,densely dashed,thick,domain=0.7:9.3,samples=60]
  {8.666667+1.466667*x+0.863439*sqrt(0.111111+(x-5)^2/60)};
\addplot[layerindigo,densely dashed,thick,domain=0.7:9.3,samples=60]
  {8.666667+1.466667*x-0.863439*sqrt(0.111111+(x-5)^2/60)};
\addplot[only marks,mark=*,mark size=1.8pt,layerindigo] coordinates {
  (1,10) (2,12) (3,13) (4,14) (5,16) (6,18) (7,19) (8,20) (9,22)};
\legend{batas prediksi, , garis pas, batas rata-rata, }
\end{axis}
\end{tikzpicture}
\caption{Selang kepercayaan $95$ persen untuk rata-rata lembur (garis putus-putus) dan
selang prediksi $95$ persen untuk satu hari baru (garis tebal), keduanya di sekitar
garis pas $\hat y=8{,}666667+1{,}466667x$. Pada $x=8$ selang rata-rata
$(19{,}9588;\ 20{,}8412)$ selebar $0{,}882$, sedangkan selang prediksi
$(19{,}4304;\ 21{,}3696)$ selebar $1{,}939$. Kedua pita melebar di ujung rentang,
tempat garis kurang ditopang data; nisbah lebar keduanya $2{,}198$ tidak bergantung
pada $n$.}
\label{fig:dua-selang}
\end{figure}

\section{Diagnostik residual}
\tingkat{menengah}

Semua selang di atas bersyarat pada model yang benar. Karena itu langkah pertama sesudah
memasang garis bukan menghitung nilai-p, melainkan memeriksa sisaan. Kalimat
sehari-harinya: kita melihat apakah yang tertinggal dari garis masih berpola; sisaan
yang berpola berarti garisnya salah bentuk, bukan kebetulan.

\index{residual}
Empat pola yang perlu dikenali: lengkungan terhadap $x$, tanda hubungannya tidak linear;
corong yang melebar bersama $x$, tanda (A3) gagal; satu titik jauh sendirian, tersangka
salah catat; dan pola menurut urutan waktu, tanda pengamatan tidak bebas.

Untuk data kafe, sisaannya berurut
\[
-0{,}1333,\ 0{,}4,\ -0{,}0667,\ -0{,}5333,\ 0,\ 0{,}5333,\ 0{,}0667,\ -0{,}4,\ 0{,}1333 .
\]
Tidak ada lengkungan yang mencolok, dan jumlahnya nol seperti seharusnya. Yang menarik
adalah ukurannya: sisaan terbesar $0{,}5333$ jam, sekitar $32$ menit, pada hari keempat
dan keenam, keduanya di dekat pusat data. Titik-titik itu justru \emph{tidak}
berpengaruh besar, karena leverage-nya hanya $0{,}1278$.

Leverage titik $i$ pada regresi sederhana punya bentuk tertutup
\[
h_{ii}=\frac1n+\frac{(x_i-\bar x)^2}{S_{xx}},
\]
sehingga pada data ini $h_{11}=h_{99}=1/9+16/60=0{,}3778$, sedangkan $h_{55}=0{,}1111$.
Titik ujung memegang $3{,}4$ kali bobot titik tengah, dan itu bukan pilihan perancang,
melainkan akibat geometri: garis yang diputar di sekitar pusat data bergerak paling
jauh di ujung. Pengaruh sebenarnya menggabungkan besarnya sisaan dengan leverage,
dan ukuran yang lazim adalah jarak Cook
\[
D_i=\frac{e_i^2}{p\,s^2}\cdot\frac{h_{ii}}{(1-h_{ii})^2},
\]
yang di sini bernilai terbesar $0{,}2870$ pada hari kedua dan kedelapan, keduanya
sisaan $0{,}4$ dengan leverage $0{,}2611$. Nilai di bawah satu umumnya dianggap tidak
mencemaskan, tetapi ambang itu konvensi, bukan hukum.

Satu peringatan yang jarang ditulis: dengan $n=9$ diagnostik hampir tidak berdaya.
Pola pada sembilan titik mudah berasal dari kebetulan, dan pola yang sungguh ada mudah
tidak terlihat. Diagnostik memberi tuduhan, bukan putusan.

\begin{bacakeras}
Dua suara.

``$R$ kuadratnya nol koma sembilan sembilan. Berarti modelnya hampir sempurna.''

``Sempurna untuk apa? Sisaannya masih dua puluh dua menit.''

``Tetapi sembilan puluh sembilan persen goyangannya sudah dijelaskan.''

``Dijelaskan oleh apa? Oleh pesanan. Bukan oleh dapur yang tiba-tiba lambat.''

``Jadi $R$ kuadrat tinggi bukan jaminan.''

``$R$ kuadrat tinggi hanya berarti garisnya rapat ke titik yang ada. Ia tidak tahu
apa pun tentang hari yang belum datang.''

Ucapkan sekali lagi, pelan: rapat ke titik lama bukan janji tentang hari baru. Ulangi
sampai terasa wajar di mulut: rapat ke titik lama, bukan janji.
\end{bacakeras}

\begin{resep}
\begin{enumerate}
  \item Gambar $y$ terhadap $x$ sebelum menghitung apa pun, dan tanyakan apakah garis
        lurus masuk akal.
  \item Pasang garis dengan kuadrat terkecil, lalu tulis persamaannya lengkap dengan
        satuan setiap koefisien.
  \item Periksa sisaan terhadap $x$ dan terhadap urutan waktu; cari lengkungan,
        corong, dan titik menyendiri.
  \item Laporkan $s$ bersama $R^2$, lalu selang untuk koefisien yang menjadi
        kesimpulan utama.
  \item Bila tujuannya meramalkan satu kejadian, pakai selang prediksi; bila tujuannya
        rata-rata, pakai selang rata-rata.
\end{enumerate}
Peringatan urutan: langkah 1 tidak boleh ditukar dengan langkah 2. Menggambar sesudah
memasang garis membuat mata mencari pembenaran, bukan mencari kesalahan.
\end{resep}

\section{Ketika jawabannya ya atau tidak}
\tingkat{lanjut}

Sampai sini $y$ berupa bilangan yang dapat diukur. Banyak pertanyaan penting berbentuk
biner: pelanggan kembali atau tidak, transaksi curang atau bukan, pasien sembuh atau
tidak. Masih menggoda memasang garis lurus biasa pada $0$ dan $1$, lalu membacanya
sebagai peluang. Dua hal salah sekaligus, dan keduanya dapat dihitung.

Data berikut: $26$ pelanggan, jarak rumah dari kedai $x$ dalam kilometer, dan $y=1$ bila
mereka kembali dalam tiga puluh hari. Tercatat $15$ dari $26$ kembali. Regresi linear
pada $y$ biner, yang disebut \emph{model peluang linear}, memberi
$\hat y=1{,}346154-0{,}096154x$. Pada $x=2$ ramalannya $1{,}153846$: peluang lebih besar
dari satu, dan itu bukan kesalahan pembulatan melainkan cacat bentuk. Pada $x=20$,
masih di luar rentang data tetapi hanya tujuh kilometer di luar, ramalannya $-0{,}576923$.

Kerusakan kedua lebih halus. Untuk $y$ biner, variansi sisaan sebenarnya $p(1-p)$ dan
karenanya bergantung pada $x$: pada $x=2$ nilainya $0{,}007225$, pada $x=14$ nilainya
$0{,}028890$. Model peluang linear mengandaikan satu ragam untuk semua titik, yaitu
$s^2=0{,}124199$ pada data ini, sehingga galat bakunya salah di tempat yang paling
penting.
Kuadrat terkecil biasa, yang bekerja baik ketika ragamnya tetap, di sini memakai
asumsi yang jelas tidak berlaku.

\index{regresi logistik}
\begin{definisi}[Model logistik]
Untuk $p_i=P(y_i=1\mid x_i)$,
\[
\log\frac{p_i}{1-p_i}=a+b\,x_i,
\qquad\text{yaitu}\qquad
p_i=\frac{1}{1+e^{-(a+bx_i)}} .
\]
Parameter ditaksir dengan kemungkinan maksimum, bukan dengan kuadrat terkecil, sehingga
tidak ada bentuk tertutup dan perhitungannya memakai iterasi.
\end{definisi}

\kamus{rasio peluang}{odds ratio}{berapa kali peluang berlipat untuk satu satuan
tambahan, bukan berapa poin peluangnya bertambah}

Untuk data itu, kemungkinan maksimum memberi $\hat a=6{,}315432$ dan $\hat b=-0{,}699945$
dengan galat baku $2{,}383413$ dan $0{,}258594$. Rasio peluangnya
$e^{-0{,}699945}=0{,}496613$: setiap kilometer tambahan memangkas peluang kembali
menjadi sekitar separuh, dalam arti nisbah, bukan selisih poin. Taksiran peluangnya
$0{,}992722$ pada $x=2$ dan $0{,}029777$ pada $x=14$, keduanya sah sebagai peluang.
Selang $95$ persen untuk $b$ adalah $(-1{,}233657;\ -0{,}166233)$; seluruhnya negatif,
jadi arahnya jelas meskipun besarnya tidak.

Satu kesalahan tafsir yang mahal: koefisien logistik bukan efek marginal pada peluang.
Nilai $b=-0{,}699945$ bekerja pada skala log nisbah, dan efeknya pada peluang paling
besar di sekitar $p=0{,}5$ serta paling kecil di ujung. Laporan yang baik menyebutkan
peluang pada beberapa nilai $x$ yang bermakna, bukan hanya koefisiennya.

\section{Korelasi bukan kausalitas}
\tingkat{dasar}

Sekarang bagian yang paling sering dilanggar di laporan yang tampak rapi. Dua kelompok
karyawan mengikuti pelatihan dengan lama berbeda, dan dicatat banyak kesalahan per
seratus transaksi.

\begin{center}
\small
\begin{tabular}{l*{6}{r}}
\hline
shift pagi  & jam latihan & 1 & 2 & 3 & 4 & 5 \\
            & kesalahan   & 4,0 & 3,6 & 3,2 & 2,9 & 2,5 \\
\hline
shift malam & jam latihan & 5 & 6 & 7 & 8 & 9 \\
            & kesalahan   & 8,0 & 7,5 & 7,0 & 6,4 & 6,0 \\
\hline
\end{tabular}
\end{center}

Dalam shift pagi, kemiringannya $-0{,}36$ kesalahan per jam latihan. Dalam shift malam,
$-0{,}474286$. Keduanya negatif: latihan lebih lama berpasangan dengan kesalahan lebih
sedikit. Tetapi bila kedua kelompok digabung, kemiringannya $+0{,}359036$ dengan
$R^2=0{,}226121$. Tanda berbalik, dan tidak ada satu pun angka yang salah.

Sebabnya bukan sihir, melainkan penjumlahan dua sumber:
\[
b_{\text{gab}}=\frac{\sum_g S_{xy,g}+\sum_g n_g(\bar x_g-\bar x)(\bar y_g-\bar y)}
{\sum_g S_{xx,g}+\sum_g n_g(\bar x_g-\bar x)^2}.
\]
Kelompok pagi menyumbang $S_{xy}=-6{,}3$ dan $S_{xx}=17{,}5$; kelompok malam
$S_{xy}=-8{,}3$ dan $S_{xx}=17{,}5$; suku antarkelompok
$6\cdot(-2)(-1{,}85)+6\cdot(2)(1{,}85)=44{,}4$; penyebutnya $83$. Maka
$b_{\text{gab}}=(-6{,}3-8{,}3+44{,}4)/83=29{,}8/83=0{,}359036$. Suku antarkelompok
$44{,}4$ mengalahkan jumlah suku dalam kelompok $-14{,}6$, dan itulah seluruh rahasianya.

Tafsirannya: jam latihan tidak diberikan secara acak. Karyawan shift malam, yang
kesalahannya lebih banyak karena jam kerja dan jenis transaksi, juga menerima latihan
lebih banyak, sehingga yang terukur sebagai efek latihan sebagian besar adalah
perbedaan shift. Gambar~\ref{fig:simpson} menaruh semuanya pada satu bidang.

\begin{figure}[htbp]
\centering
\begin{tikzpicture}
\begin{axis}[
  width=11.4cm, height=6.2cm,
  xmin=0.5, xmax=10.5, ymin=1.6, ymax=8.9,
  axis lines=left,
  xtick={1,2,3,4,5,6,7,8,9,10},
  tick label style={font=\scriptsize},
  xlabel={jam latihan}, ylabel={kesalahan per 100 transaksi},
  xlabel style={font=\small}, ylabel style={font=\small},
  legend style={font=\scriptsize,draw=layergrey,at={(0.98,0.98)},
    anchor=north east,row sep=0.5pt},
]
\addplot[only marks,mark=*,mark size=2pt,layerteal] coordinates {
  (1,4.0) (2,3.6) (3,3.2) (4,2.9) (5,2.5)};
\addplot[only marks,mark=square*,mark size=1.9pt,layerplum] coordinates {
  (6,7.5) (7,7.0) (8,6.4) (9,6.0) (10,5.7)};
\addplot[layerteal,very thick,domain=0.8:5.4,samples=2] {4.326667-0.36*x};
\addplot[layerplum,very thick,domain=5.6:10.2,samples=2]
  {10.32381-0.474286*x};
\addplot[accent,thick,densely dashed,domain=0.8:10.2,samples=2]
  {2.941968+0.359036*x};
\legend{pagi, malam, gabungan}
\end{axis}
\end{tikzpicture}
\caption{Paradoks Simpson pada data pelatihan. Di dalam setiap shift, garisnya menurun:
$-0{,}36$ kesalahan per jam untuk pagi dan $-0{,}4743$ untuk malam. Garis gabungan,
yang dihitung dari kedua belas titik tanpa memandang shift, justru menanjak dengan
kemiringan $+0{,}3590$. Perubahan tanda itu berasal dari selisih antar kelompok, bukan
dari satu titik yang aneh; suku antarkelompok $44{,}4$ mengalahkan suku dalam kelompok
$-14{,}6$.}
\label{fig:simpson}
\end{figure}

\begin{untukanak}
Dua kelas mengikuti lomba membaca. Kelas A membaca buku tipis, banyak halaman, dan
nilainya bagus. Kelas B membaca buku tebal, lebih sedikit halaman, dan nilainya jelek.
Kalau semua murid digabung, terlihat seolah-olah makin sedikit halaman yang dibaca,
makin tinggi nilainya. Padahal di dalam setiap kelas, makin banyak halaman, makin tinggi
nilainya.

Yang salah bukan angkanya, melainkan pencampurannya. Kelas B memang diberi buku yang
lebih tebal, dan kelas B juga memang lebih lemah membaca. Dua hal itu bercampur menjadi
satu garis yang menipu.

Di bab ini kita akan tahu mengapa mencampur dua kelompok dapat membalik arah kesimpulan,
mengapa pertanyaan ``apakah latihan mengurangi kesalahan'' tidak dapat dijawab hanya
dengan menggambar semua titik, dan mengapa orang yang mengubah kebijakan harus tahu
kelompok mana yang sebenarnya ia bandingkan.
\end{untukanak}

\begin{tantangan}
Susun sebuah tabel dengan dua kelompok, masing-masing lima titik, sedemikian sehingga
kedua kemiringan dalam kelompok positif tetapi kemiringan gabungan negatif, dan
selisih rata-rata $y$ antarkelompok paling kecil yang masih memungkinkan. Berapa
selisih terkecil itu, dan mengapa tidak bisa lebih kecil?

Petunjuk 1: tanda kemiringan gabungan ditentukan oleh suku antarkelompok. Buat
kelompok pertama memiliki $x$ kecil dan $y$ kecil, kelompok kedua $x$ besar dan $y$
besar, lalu balik arah di dalam masing-masing kelompok.

Petunjuk 2: dengan lima titik dan kemiringan dalam kelompok $+1$ pada
$S_{xx}=10$, setiap kelompok menyumbang $S_{xy}=+10$ dan suku antarkelompok harus
lebih kecil dari $-20$. Karena $n_g(\bar x_g-\bar x)(\bar y_g-\bar y)$ bernilai
negatif hanya bila salah satu selisih rata-rata berlawanan tanda dengan yang lain,
selisih rata-rata harus melewati titik di mana kedua suku itu saling meniadakan.

Kalau mentok, kembali ke bagian Korelasi bukan kausalitas dan hitung ulang
$29{,}8/83$ langkah demi langkah.
\end{tantangan}

\section{Pelaporan yang jujur}
\tingkat{menengah}

Sebuah analisis regresi dapat dikerjakan dengan benar dan dilaporkan sehingga menyesatkan
tanpa satu pun angka dipalsukan. Empat kebiasaan menutup sebagian besar celah itu.

Pertama, sebutkan berapa pengamatan dan berapa parameter, karena pembaca memerlukan
keduanya untuk menghitung derajat kebebasan dan menilai apakah $s$ dapat dipercaya.
Pada data kafe, $n=9$ dan $p=2$, jadi $s$ dibangun dari tujuh derajat kebebasan;
menyembunyikan angka itu membuat $s$ tampak lebih kokoh daripada kenyataannya.

Kedua, pisahkan data yang dipakai memilih model dari data yang dipakai menilai.
Memilih bentuk model, memilih peubah, dan menghitung galat pada himpunan yang sama
selalu menghasilkan angka yang terlalu bagus. Pada $n=9$ pemisahan itu tidak mungkin,
dan justru karena itu kesimpulannya harus disebut lemah, bukan diperkuat dengan
$R^2=0{,}992821$.

Ketiga, sebutkan berapa banyak model yang dicoba. Bila sepuluh kemungkinan diuji dan
yang terbaik dilaporkan, peluang sekurangnya satu di antaranya tampak berarti adalah
$1-0{,}95^{10}=0{,}401263$ meskipun tidak satu pun efeknya nyata. Angka itu bukan
peluang salah, melainkan peluang terpilihnya sesuatu yang kebetulan.

Keempat, laporkan besaran dengan satuan dan selang, bukan hanya tanda bintang. Kalimat
``$\hat b=0{,}687$ ($95\%$: $0{,}381$ sampai $0{,}993$ juta per karyawan)'' dapat
diperiksa; kalimat ``berpengaruh signifikan'' tidak. Bila kesimpulannya tentang kausalitas,
sebutkan mengapa: pengacakan, atau penyesuaian terhadap pencampur yang disebut namanya.

\section{Contoh hitung dari ujung ke ujung}
\tingkat{menengah}

Kita ulang seluruh rantai pada data kafe, dengan angka antara ditulis, supaya setiap
langkah dapat diperiksa dengan kalkulator.

Rata-ratanya $\bar x=5$ dan $\bar y=16$. Simpangan kuadratnya
$S_{xx}=\sum(x_i-5)^2=60$ dan $S_{yy}=130$; hasil kalinya
$S_{xy}=\sum(x_i-5)(y_i-16)=808-9\cdot5\cdot16=88$. Maka
$\hat b=88/60=1{,}466667$ dan $\hat a=16-1{,}466667\cdot5=8{,}666667$.

Ramalan pada sembilan titik itu $10{,}1333$; $11{,}6$; $13{,}0667$; $14{,}5333$;
$16$; $17{,}4667$; $18{,}9333$; $20{,}4$; $21{,}8667$. Sisaannya
$-0{,}1333$; $0{,}4$; $-0{,}0667$; $-0{,}5333$; $0$; $0{,}5333$; $0{,}0667$;
$-0{,}4$; $0{,}1333$, dan jumlah kuadratnya
$\mathrm{SSE}=0{,}933333$. Dengan $n-p=7$, $s^2=0{,}133333$ dan $s=0{,}365148$ jam.

Galat baku kemiringan $s/\sqrt{S_{xx}}=\sqrt{0{,}133333/60}=\sqrt{0{,}00222222}
=0{,}047140$. Statistik $t=1{,}466667/0{,}047140=31{,}1127$, dan
$t_{7,0{,}975}=2{,}364624$, sehingga setengah lebar selangnya $0{,}111470$ dan
selangnya $(1{,}355197;\ 1{,}578136)$.

\kertas{Ulangi perhitungan yang sama untuk satu perubahan data. Hari kesembilan
tercatat $22$ jam lembur; andaikan catatan yang benar adalah $26$ jam. Hitung ulang
$\hat b$, $s$, $\operatorname{se}(\hat b)$, dan selang $95$ persennya, lalu bandingkan
dengan tabel di atas. Jawaban yang benar: $\hat b=1{,}733333$, $s=1{,}306880$,
$\operatorname{se}(\hat b)=0{,}168718$, selang $(1{,}3344;\ 2{,}1323)$. Lebarnya
$0{,}798$ melawan $0{,}223$, yaitu $3{,}6$ kali. Jelaskan mengapa satu titik di ujung
dapat melakukan itu, padahal ia hanya satu dari sembilan.}

\section{Lab komputasi}
\tingkat{lanjut}

Semua angka di bab ini dapat dihitung ulang dalam satu berkas pendek. Program berikut
memakai $X^{\top}X$ secara langsung, jadi ia sekaligus memeriksa aljabar di bagian
kedua.

\begin{lab}
\begin{verbatim}
import numpy as np
from math import sqrt
from scipy import stats

x = np.arange(1, 10, dtype=float)
y = np.array([10, 12, 13, 14, 16, 18, 19, 20, 22], dtype=float)
n = len(y)
X = np.column_stack([np.ones(n), x])

beta = np.linalg.solve(X.T @ X, X.T @ y)
fit = X @ beta
e = y - fit
s2 = (e ** 2).sum() / (n - 2)
sst = ((y - y.mean()) ** 2).sum()

print("X'X =", (X.T @ X).tolist(), " X'y =", (X.T @ y).tolist())
print("beta = %.6f %.6f" % tuple(beta))
sse = (e ** 2).sum()
print("SSE = %.6f  R2 = %.6f  s = %.6f"
      % (sse, 1 - sse / sst, sqrt(s2)))
print("residu =", np.round(e, 4).tolist())

cov = s2 * np.linalg.inv(X.T @ X)
se = np.sqrt(np.diag(cov))
tq = stats.t.ppf(0.975, n - 2)
print("se = %.6f %.6f  t = %.4f %.4f"
      % (se[0], se[1], beta[0] / se[0], beta[1] / se[1]))
lo, hi = beta[1] - tq * se[1], beta[1] + tq * se[1]
print("IK95 kemiringan = (%.6f, %.6f)" % (lo, hi))

for x0 in (5.0, 8.0):
    c = np.array([1.0, x0])
    y0 = c @ beta
    v = s2 * (c @ np.linalg.inv(X.T @ X) @ c)
    print("x0=%.0f  yhat=%.4f  rata2 +/-%.4f  prediksi +/-%.4f"
          % (x0, y0, tq * sqrt(v), tq * sqrt(v + s2)))

h = np.einsum("ij,jk,ik->i", X, np.linalg.inv(X.T @ X), X)
cook = e ** 2 / (2 * s2) * h / (1 - h) ** 2
print("h_ii =", np.round(h, 4).tolist())
print("D_Cook =", np.round(cook, 4).tolist())

x1 = np.array([6, 7, 7, 8, 8, 9, 9, 10, 10, 11, 12, 12], dtype=float)
x2 = np.array([18, 22, 20, 26, 24, 28, 26, 32, 30, 34, 36, 38], dtype=float)
yb = np.array([11.2, 12.4, 12.0, 14.1, 13.6, 15.2, 14.8, 17.0, 16.4, 18.1,
               19.4, 19.9])
Xb = np.column_stack([np.ones(12), x1, x2])
bb = np.linalg.solve(Xb.T @ Xb, Xb.T @ yb)
eb = yb - Xb @ bb
s2b = (eb ** 2).sum() / 9
seb = np.sqrt(np.diag(s2b * np.linalg.inv(Xb.T @ Xb)))
r12 = np.corrcoef(x1, x2)[0, 1]
print("berganda beta = %.4f %.4f %.4f  s = %.4f"
      % (bb[0], bb[1], bb[2], sqrt(s2b)))
print("berganda se   = %.4f %.4f %.4f" % tuple(seb))
print("r12 = %.4f  VIF = %.2f  akar VIF = %.2f"
      % (r12, 1 / (1 - r12 ** 2), 1 / sqrt(1 - r12 ** 2)))


def ols(xx, yy):
    A = np.column_stack([np.ones(len(xx)), xx])
    return np.linalg.solve(A.T @ A, A.T @ yy)[1]


xp = np.array([1, 2, 3, 4, 5, 6], dtype=float)
yp = np.array([4.0, 3.6, 3.2, 2.9, 2.5, 2.2])
xm = np.array([5, 6, 7, 8, 9, 10], dtype=float)
ym = np.array([8.0, 7.5, 7.0, 6.4, 6.0, 5.7])
print("Simpson: pagi %.4f  malam %.4f  gabung %.4f"
      % (ols(xp, yp), ols(xm, ym), ols(np.r_[xp, xm], np.r_[yp, ym])))

xd = np.array([2, 2, 3, 3, 4, 4, 5, 5, 6, 6, 7, 7, 8, 8, 9, 9, 10, 10,
               11, 11, 12, 12, 13, 13, 14, 14], dtype=float)
yd = np.array([1, 1, 1, 1, 1, 1, 1, 1, 1, 1, 1, 0, 1, 0, 1, 0, 1, 0, 0,
               1, 0, 0, 0, 0, 0, 0], dtype=float)
A = np.column_stack([np.ones(len(yd)), xd])
lpm = np.linalg.solve(A.T @ A, A.T @ yd)
print("LPM: %.4f %+.4f  prediksi di x=2: %.4f  di x=20: %.4f"
      % (lpm[0], lpm[1], lpm @ [1, 2], lpm @ [1, 20]))


def nll(th):
    z = A @ th
    return float(np.sum(np.log1p(np.exp(-z)) + (1 - yd) * z))


g = np.array([0.0, 0.0])
for _ in range(40):
    p = 1 / (1 + np.exp(-(A @ g)))
    grad = A.T @ (p - yd)
    Hm = A.T @ (A * (p * (1 - p))[:, None])
    g = g - np.linalg.solve(Hm, grad)
se_g = np.sqrt(np.diag(np.linalg.inv(Hm)))
p = 1 / (1 + np.exp(-(A @ g)))
print("logit: %.4f %+.4f  se %.4f %.4f" % (g[0], g[1], se_g[0], se_g[1]))
print("devians = %.4f  p(x=2) = %.4f  p(x=14) = %.4f"
      % (2 * nll(g), p[0], p[-1]))
print("rasio peluang per km = %.4f" % np.exp(g[1]))
\end{verbatim}
\end{lab}

Keluarannya, apa adanya:

\begin{verbatim}
X'X = [[9.0, 45.0], [45.0, 285.0]]  X'y = [144.0, 808.0]
beta = 8.666667 1.466667
SSE = 0.933333  R2 = 0.992821  s = 0.365148
residu = [-0.1333, 0.4, -0.0667, -0.5333, 0.0, 0.5333, 0.0667, -0.4, 0.1333]
se = 0.265274 0.047140  t = 32.6706 31.1127
IK95 kemiringan = (1.355197, 1.578136)
x0=5  yhat=16.0000  rata2 +/-0.2878  prediksi +/-0.9101
x0=8  yhat=20.4000  rata2 +/-0.4412  prediksi +/-0.9696
h_ii = [0.3778, 0.2611, 0.1778, 0.1278, 0.1111, 0.1278, 0.1778, 0.2611, 0.3778]
D_Cook = [0.0651, 0.287, 0.0044, 0.1792, 0.0, 0.1792, 0.0044, 0.287, 0.0651]
berganda beta = 2.4430 0.6869 0.2393  s = 0.1436
berganda se   = 0.2076 0.1354 0.0421
r12 = 0.9868  VIF = 38.13  akar VIF = 6.18
Simpson: pagi -0.3600  malam -0.4743  gabung 0.3590
LPM: 1.3462 -0.0962  prediksi di x=2: 1.1538  di x=20: -0.5769
logit: 6.3154 -0.6999  se 2.3834 0.2586
devians = 17.8843  p(x=2) = 0.9927  p(x=14) = 0.0298
rasio peluang per km = 0.4966
\end{verbatim}

Dua hal perlu dibaca dari keluaran itu. Pertama, $\mathrm{X}^{\top}\mathrm{X}$ dan
$\mathrm{X}^{\top}y$ tercetak persis seperti dihitung tangan, dan $\hat\beta$ keluar
dari penyelesaian sistem; kesepakatan semacam ini bukan bukti kebenaran teori, tetapi ia
menyingkirkan salah hitung, yang paling sering menjadi penyebab. Kedua, tiga bagian
terakhir memperlihatkan mengapa tafsiran lebih rapuh daripada hitungan: $0{,}6869$
berbanding $1{,}4460$ pada regresi tunggal, VIF $38{,}13$, dan kemiringan gabungan yang
tandanya berbalik. Tidak satu pun dari ketiganya kesalahan program.

\begin{untukdosen}
Urutan mengajar yang saya sarankan: mulai dari tabel sembilan hari, minta kelas menebak
kemiringannya dengan dua cara, yaitu selisih ujung dan ``garis yang paling dekat''.
Selisih kedua jawaban kecil, dan justru itu pintu masuknya: pertanyaan sebenarnya bukan
angkanya, melainkan goyangannya. Baru sesudah kelas gelisah, masuk ke persamaan normal
dan tunjukkan bahwa $\sum x_ie_i=0$ adalah pernyataan geometris, bukan trik aljabar.

Dua miskonsepsi yang paling sering muncul. Pertama, mahasiswa membaca $R^2$ sebagai
``peluang ramalan benar'', lalu terkejut ketika selang prediksi selebar dua jam pada
$R^2=0{,}99$; gejalanya tampak ketika mereka tidak dapat menjelaskan mengapa $s$ tidak
ikut mengecil bersama $R^2$. Kedua, mereka menafsirkan koefisien berganda seperti
koefisien tunggal, padahal pada data dua belas cabang angkanya berubah dari $1{,}446$
menjadi $0{,}687$ hanya karena penjelas lain masuk.

Aktivitas papan tulis yang berhasil: suruh dua kelompok menghitung kemiringan pada data
pelatihan, satu kelompok memakai semua dua belas titik, satu kelompok memisahkan shift.
Tulis kedua hasil berdampingan, lalu minta masing-masing mempertahankannya. Tanda bahwa
kelas siap lanjut: mereka dapat menyebut satu keadaan yang membuat koefisien berganda
berubah tanda, dan dapat menjelaskan mengapa ramalannya tetap baik.
\end{untukdosen}

\section{Ringkasan}
\tingkat{dasar}

\begin{ringkasan}
\begin{itemize}
  \item Regresi memasang garis yang meminimumkan jumlah kuadrat sisaan. Menurut
        persamaan normal sisaannya tegak lurus setiap kolom $X$; pada data kafe
        $\hat b=88/60=1{,}466667$ dan $\hat a=8{,}666667$.
  \item Koefisien hanya bermakna bersama satuannya: $1{,}466667$ jam per puluhan pesanan
        sama dengan $8{,}8$ menit per pesanan. Mengubah skala $x$ mengubah koefisien dan
        galat bakunya, tidak mengubah $t$ maupun $R^2$.
  \item $R^2$ nisbi, $s$ mutlak. Pada data kafe $R^2=0{,}992821$ tetapi
        $s=0{,}365148$ jam, yaitu $21{,}9$ menit.
  \item Di bawah (A1)--(A4) Gauss--Markov menjamin kuadrat terkecil adalah BLUE dengan
        $\operatorname{Var}(\hat\beta)=\sigma^2(X^{\top}X)^{-1}$; jaminan itu gugur bila
        ragam galat tidak tetap atau penjelasnya kolinear.
  \item Kolinearitas tidak merusak ramalan tetapi merusak tafsiran: pada dua belas cabang,
        $r_{12}=0{,}986801$ memberi $\mathrm{VIF}=38{,}13$ dan galat baku $6{,}18$ kali
        lebih besar, sementara $R^2$ tetap $0{,}997951$.
  \item Selang prediksi selalu lebih lebar daripada selang rata-rata; pada $x_0=8$
        nisbahnya $2{,}198$ dan tidak mengecil bersama $n$. Tanda kemiringan dapat
        berbalik ketika dua kelompok dicampur: $-0{,}36$ dan $-0{,}474286$ di dalam
        kelompok, $+0{,}359036$ setelah digabung.
\end{itemize}
\end{ringkasan}

\section{Latihan}
\tingkat{menengah}

\begin{latihan}
\begin{enumerate}
  \item Nilai klaim ini: ``Karena $R^2=0{,}99$, garis ini dapat meramalkan satu hari
        dengan galat paling banyak dua puluh menit.'' Benar, salah, atau belum cukup
        bukti? Sebutkan bagian yang benar, bagian yang keliru, dan angka tambahan yang
        menyelesaikan penilaian itu.
  \item Buktikan bahwa dengan kolom satu pada $X$ berlaku $\sum_i e_i=0$ dan
        $\sum_i x_ie_i=0$, lalu jelaskan mengapa keduanya gagal bila garis dipaksa lewat
        titik asal.
  \item Buktikan $\operatorname{Var}(\hat b)=\sigma^2/S_{xx}$ dari
        $\hat b=\sum_i(x_i-\bar x)y_i/S_{xx}$ tanpa rumus matriks.
  \item Hitung $\hat a$, $\hat b$, $s$, $\operatorname{se}(\hat b)$, dan selang $95$
        persen untuk $(1,2)$, $(2,2{,}5)$, $(3,4)$, $(4,3{,}5)$, $(5,5)$, lalu periksa
        dengan Python. Mengapa selangnya jauh lebih lebar daripada selang data kafe?
  \item Hitung selang prediksi $95$ persen untuk $x_0=2$ dan $x_0=9$ pada data kafe.
        Mengapa lebarnya berbeda padahal $n$ sama, dan mengapa selang itu tidak pernah
        menyempit ke nol ketika $n$ diperbesar?
  \item Tunjukkan bahwa menambah penjelas tidak dapat menaikkan SSE, lalu jelaskan
        mengapa $R^2_{\text{adj}}$ dapat turun. Hitung $R^2_{\text{adj}}$ untuk data kafe
        bila $p$ naik dari $2$ menjadi $3$ pada $R^2$ yang sama.
  \item Untuk $S_{11}=50$, $S_{22}=200$, $r_{12}=0{,}9$, dan $\sigma^2=2$, hitung VIF
        serta galat baku $\hat b_1$ dengan dan tanpa $x_2$ dalam model.
  \item Susun tabel dua kelompok dengan kemiringan dalam kelompok $-0{,}2$ dan kemiringan
        gabungan $+0{,}4$, lalu periksa hasilnya dengan rumus dekomposisi.
  \item Tulis ulang kalimat ini sehingga jujur: ``model terbaik dari dua belas model yang
        kami coba memberi $p=0{,}03$, jadi efeknya nyata''.
\end{enumerate}

Petunjuk untuk tiga butir tersulit. Butir 3: tulis $\hat b-\beta=\sum_i(x_i-\bar x)
\varepsilon_i/S_{xx}$, ambil variansi, pakai ketakbebasan dan
$\operatorname{Var}(\varepsilon_i)=\sigma^2$. Butir 6: bandingkan ruang kolom kedua model,
lalu ingat pembagi $n-p$ yang mengecil. Butir 7: $\mathrm{VIF}=1/(1-0{,}81)=5{,}263158$,
jadi galat bakunya $\sqrt{5{,}263158}$ kali lebih besar.
\end{latihan}

\begin{untukpeneliti}
Yang benar-benar terbukti di sini: persamaan normal sebagai syarat perlu dan cukup bagi
minimum beserta ketunggalannya saat $X$ berpangkat kolom penuh; sifat matriks topi dan
$\operatorname{tr}H=p$; ketakbiasan serta
$\operatorname{Var}(\hat\beta)=\sigma^2(X^{\top}X)^{-1}$; ketidaksamaan Gauss--Markov
dalam kelas penaksir linear tak bias; dekomposisi $\mathrm{SST}=\mathrm{SSR}+\mathrm{SSE}$
bila ada kolom satu; ketakmonotonan $R^2$ terhadap penambahan kolom; sebaran $t$ untuk
nisbah koefisien bila galat normal; dan variansi galat prediksi.

Yang hanya populer: bahwa $R^2$ mengukur mutu model, bahwa koefisien berganda selalu
merupakan efek, dan bahwa ambang $\mathrm{VIF}>10$ atau $D_i>1$ adalah hasil turunan.
Sumber primer: Seber dan Lee, \emph{Linear Regression Analysis}, Bab 3; Greene,
\emph{Econometric Analysis}, Bab 3; Wasserman, \emph{All of Statistics}, Bab 13;
Simpson (1951) untuk pembalikan arah pada tabel kontingensi; Hosmer, Lemeshow, dan
Sturdivant, \emph{Applied Logistic Regression}, Bab 1--3.

Batas keberlakuan: semua selang bersyarat pada model yang benar, galat bebas, dan ragam
tetap. Pada data berkelompok, deret waktu, atau pengamatan berulang, ketakbebasan gagal
dan galat baku yang dilaporkan terlalu kecil. Pertanyaan terbuka: seberapa jauh selang
pascapemilihan model menyimpang dari nominalnya, dan bagaimana diagnostik pengaruh
disusun ketika $n$ jauh lebih kecil daripada $p$.
\end{untukpeneliti}

\peta{Bab 21 pada Wasserman, \emph{All of Statistics}: Bab 13 (Linear and Logistic
Regression: kuadrat terkecil, sebaran $t$ untuk koefisien, dan regresi logistik) serta
Bab 6 (Models, Statistical Inference, and Learning: tafsiran selang dan galat baku).
Geometri kuadrat terkecil sebagai proyeksi ortogonal disusun di Bab 10 buku ini, tempat
persamaan normal pertama kali muncul sebagai akibat ortogonalitas, dan dipakai lagi di
Bab 24 untuk regularisasi. Selang kepercayaan, nilai-p, dan tafsiran taraf diambil dari
Bab 20 tanpa diulang di sini. Teorema Gauss--Markov beserta syaratnya dirujuk ke Seber
dan Lee Bab 3 dan Greene Bab 3, karena Wasserman tidak membuktikannya. Paradoks Simpson
dirujuk ke makalah Simpson (1951) dan dipakai lagi di Bab 25 ketika membahas pemisahan
data dan kebocoran evaluasi, serta di Bab 26 ketika membahas klaim kausal yang harus
dibuktikan, bukan disimpulkan dari korelasi.}
